\section{Solucions i eines existents}

Les propostes revisades es poden agrupar en simuladors acadèmics, biblioteques
generals de simulació i cues operatives. GridSim modela recursos Grid, usuaris
i aplicacions, i permet estudiar algorismes d'assignació sota recursos
heterogenis \parencite{buyya2002gridsim}. CloudSim trasllada aquest enfocament
a centres de dades, màquines virtuals i polítiques de provisió cloud
\parencite{calheiros2011cloudsim}. SimGrid se centra en aplicacions distribuïdes
i ofereix abstraccions per experimentar amb algorismes de planificació. Els seus
autors destaquen que els simuladors creats per a un sol estudi dificulten
reproduir i comparar resultats \parencite{legrand2003simgrid}.

Aquestes eines demostren el valor d'intercanviar polítiques sense reconstruir
tot el model, però inclouen xarxes, recursos heterogenis, VM o aplicacions
paral·leles. SimPy ofereix una alternativa general per construir simulacions de
processos i recursos en Python \parencite{simpyTime}; proporciona mecanismes
temporals, però cal construir el model, les polítiques, la persistència i la
interfície.

Com a antecedent universitari, SimuQ permet definir sistemes de cues des d'una
interfície, executar-los amb SimPy i obtenir dades per comparar configuracions
\parencite{caballero2016simuq}. És proper per la combinació de simulació i
interacció gràfica, però representa punts de servei, no tasques backend, i és
una aplicació d'escriptori.

BullMQ i RabbitMQ pertanyen a la família operativa. BullMQ gestiona tasques i
workers sobre Redis i permet FIFO, LIFO i prioritats
\parencite{bullmqPrioritized}. RabbitMQ és un broker de missatges amb cues FIFO
i prioritats subjectes al tipus de cua i a la configuració
\parencite{rabbitmqPriority}. Cap dels dos està dissenyat per executar la mateixa
càrrega sintètica amb diverses polítiques i presentar una comparació neutral.

\begin{table}[htbp]
  \centering
  \caption{Comparació resumida de solucions relacionades}
  \label{tab:eines}
  \begin{tabularx}{\textwidth}{@{}p{1.9cm}p{2.4cm}X X@{}}
    \toprule
    Proposta & Tipus & Aportació principal & Diferència respecte del TFM \\
    \midrule
    GridSim & Simulador & Recursos Grid i polítiques extensibles & Domini distribuït, terminis i costos \\
    CloudSim & Simulador & Infraestructura cloud i provisió & Modela centres de dades, VM i xarxes \\
    SimGrid & Marc & Planificació d'aplicacions distribuïdes & Orientació Grid/HPC i configuració tècnica \\
    SimPy & Biblioteca & Processos, esdeveniments i recursos & Cal construir l'aplicació experimental completa \\
    SimuQ & Eina acadèmica & Interfície per simular cues de servei & Escriptori i domini de punts de servei \\
    BullMQ / RabbitMQ & Cua operativa & Execució real, workers i prioritats & No comparen polítiques sota entrades equivalents \\
    \bottomrule
  \end{tabularx}
\end{table}

La comparació no identifica un competidor directe. Mostra un espai entre els
marcs de simulació potents però tècnics i les cues productives que executen una
configuració concreta. Aquest espai permet situar una plataforma web acotada,
orientada a preparar escenaris, repetir-los i interpretar els compromisos entre
polítiques.
